% Chapitre : module de besoin (demande de soins de ville), indice de couverture et transferts de tâches.
% Sources principales : v2/papers/besoin/main.tex (article complet, à jour de l'erratum P1-4 du 2026-11-15),
% v2/docs/lab/reports/2026-11-15-besoin-*.md, v2/docs/lab/derivations/2026-11-15-besoin-derivation.md,
% v2/besoin/ (code), v2/docs/lab/decisions.md (entrées « besoin » et « L23 » du 2026-11-15).
% Ce chapitre résume l'article ; les tableaux détaillés des protocoles sont dans son annexe.
\begingroup\sloppy\setlength{\emergencystretch}{3em}% coupures des chemins longs, local à ce fichier

\section{Module de besoin : demande de soins de ville et indice de couverture}\label{sec:besoin-ref}\label{bes:sec}

Le jumeau dit \emph{combien} de professionnels il y aura. Pour dire si ce sera \emph{assez}, il faut un dénominateur.
Ce chapitre résume le module qui le fournit ; l'article complet est \texttt{v2/papers/besoin/main.pdf}, ses tableaux
détaillés sont en annexe de l'article. Le code est dans \texttt{v2/besoin/}. La dérivation est
\texttt{v2/docs/lab/derivations/2026-11-15-besoin-derivation.md}, et les rapports de passage sont
\texttt{v2/docs/lab/reports/2026-11-15-besoin-*.md}. Le module ne modifie pas la projection des effectifs : il
ajoute un indice à côté.

\subsection{Le problème : on observe le recours, pas le besoin}\label{bes:probleme}
Pour une cellule $i$ (territoire, âge, sexe) et une profession $p$, le recours observé s'écrit
\[
R_{ip} = B_{ip}\,\varphi_p(A_{ip}) + D_p(A_{ip}) + \varepsilon_{ip},\qquad 0<\varphi\le 1,\quad D\ge 0,
\]
où $B$ est le besoin selon une norme déclarée, $A$ l'accessibilité de l'offre, $\varphi$ la part du besoin satisfaite
et $D$ la demande induite par l'offre~\citep{evans1974,delattre2003,leonard2009}. Avec $D=0$, le couple
$(B\,h,\ \varphi/h)$ produit le même $R$ pour toute fonction $h>0$\derive{} : le besoin n'est pas identifiable sans
restriction (indépendance conditionnelle, saturation, ou variable d'exclusion), ou sans norme externe à l'offre.

\paragraph{Ce que l'indexation sauve.} L'indice servi est
\[
C(p,T,t)=\frac{\text{offre}(p,T,t)\,/\,B(p,T,t)}{\text{offre}(p,T,2026)\,/\,B(p,T,2026)} ,
\]
qui vaut 1 en 2026. Tout biais multiplicatif constant par territoire s'y annule\derive{} : comparer 2035 à 2026 dans
un territoire est défendable avec un estimateur simple ; comparer deux territoires en niveau ne l'est pas.

\paragraph{Ce que l'indice mesure.} Le recours ouvert est celui des soins de ville remboursés, hôpital public exclu.
L'indice mesure donc la couverture de la \emph{demande de soins de ville adressée au secteur privé}, pas le besoin de
toute la population. La carte et l'assistant l'appellent \enquote{Couverture de la demande de ville} ; l'API porte
l'attribut \texttt{grandeur}.

\subsection{Six estimateurs concurrents}\label{bes:estimateurs}
Tous passent par l'interface \texttt{v2/besoin/interface.py} (\texttt{fit}, \texttt{predict}, \texttt{assumptions},
\texttt{norme}, \texttt{grandeur}), conformément à la règle du projet : on compare des estimateurs, on ne présente
pas un chiffre unique.

\begin{center}\small
\begin{tabularx}{\linewidth}{@{}l X X@{}}
\toprule
& \textbf{Forme} & \textbf{Restriction qui l'identifie} \\
\midrule
M0 & $B_0=\bar r_p\,\mathrm{pop}(T,t)$ & besoin par habitant uniforme ; c'est la densité de la carte \\
M1 & $B_1=\sum_{a,s}\mathrm{pop}(T,t,a,s)\,r(p,a,s)$ & besoin par âge et sexe = recours national par âge et sexe \\
M2a & recours régressé sur besoin et offre, évalué à l'offre moyenne nationale & neutralisation de l'offre~\citep{carrhill1994,gravelle2003} \\
M2b & $R_{\mathrm{std}}=\rho_{\max}A/(A+K)$, besoin au plateau & saturation (extension testée ici) \\
M3 & $\sum_g \mathrm{pop}\cdot\mathrm{prev}(g)\cdot n(p,g)$ & norme : prévalence $\times$ soins par patient~\citep{birch2007,mackenzie2019} \\
M4 & 2SFCA à paliers & aucune : indice d'accès qui prend le besoin de M1 à M3 en entrée~\citep{luowang2003,luo2009} \\
\bottomrule
\end{tabularx}
\end{center}

\begin{itemize}
  \item \textbf{M1} prend la forme du profil d'âge dans les poids publiés avec l'APL de la DREES (consommation DCIR
        2020--2022, 19 tranches) et le niveau dans Open DAMIR 2025. Pour les infirmiers, les 90~ans et plus pèsent
        15,5 fois la moyenne.
  \item \textbf{M3} combine la cartographie des pathologies de la CNAM (61 groupes, 2015--2024) et les dépenses par
        pathologie converties en actes ; cette combinaison est la nôtre.
  \item \textbf{M4} reprend les paliers de l'APL 2024~\citep{vergier2017} avec des temps routiers OSRM sur
        OpenStreetMap, et fige les poids d'âge à une année de base pour rester projetable. Il n'est jamais appelé
        \enquote{APL} : c'est un indice d'accessibilité maison.
\end{itemize}

\subsection{Sept protocoles gelés et horodatés}\label{bes:validation}
Chaque protocole est écrit et gelé avant tout calcul. L'empreinte SHA-256 du manifeste est horodatée par deux
autorités RFC~3161 (FreeTSA, DigiCert) et soumise à OpenTimestamps, parce que l'historique git ne prouve pas
l'antériorité d'un gel (chapitre~\ref{sec:validation-ref}). Aucun seuil n'a été changé après un résultat ; tout écart
au protocole est écrit au registre avant le calcul.

\begin{center}\small
\begin{tabularx}{\linewidth}{@{}l X X@{}}
\toprule
\textbf{Test} & \textbf{Question} & \textbf{Verdict} \\
\midrule
P-B1 & L'âge (M1) déplace-t-il l'indice 2050 par rapport à M0 ? & M1 retenu pour les cinq professions ; IDE et MK
$S_1=0$\,\%\mesure{} ; dentistes de justesse ($S_1=87{,}1$\,\%\mesure{} pour un seuil de 90\,\%) \\
P-B2 & Le recours plafonne-t-il quand l'accès augmente ? & M2b rejeté pour IDE, MK, CD, SF : le besoin non
satisfait n'est pas identifié en données ouvertes \\
P-B3 & La morbidité (M3) apporte-t-elle au-delà de l'âge ? & Retenue pour l'indice des seuls IDE
($S_3=83$\,\%\mesure{}) ; pouvoir explicatif hors échantillon rejeté ($\Delta R^2$ de $-0{,}03$ à
$+0{,}01$\mesure{}) \\
P-B4 & M4 reproduit-il l'APL 2024 ? & Échec : Spearman 0,71 à 0,76\mesure{} pour un seuil de 0,90, mais gain de
$+0{,}15$ à $+0{,}35$\mesure{} sur la meilleure densité ; accord départemental $\rho=0{,}98$\mesure{} (MK) \\
P-B4-PH & M4 reproduit-il l'APL V0 2013 des pharmacies ? & Échec : $\rho=0{,}72$\mesure{}, erreur médiane
33\,\%\mesure{} ; niveau moyen reproduit \\
P-B5 & Backtest 2016\,$\to$\,2025 : l'âge bat-il la persistance ? & Oui pour 3 professions sur 4 (IDE
$-8{,}07$, MK $-3{,}32$, CD $-0{,}99$ points de MAPE\mesure{}) ; non pour les SF ($+1{,}57$\mesure{}) \\
P-B6 & La vaccination en officine a-t-elle ajouté des vaccinations ? & ATT $+0{,}73$\,\%
[$-0{,}12$\,;\,$+1{,}58$]\mesure{}, placebo en échec : non attribuable, l'expansion $m$ reste supposée à 0 \\
\bottomrule
\end{tabularx}
\end{center}

Le GBM concurrent de P-B5 bat M1 au backtest, mais son gain suit la croissance de l'offre (corrélation 0,54 à 0,71,
diagnostic a posteriori) : il n'est servi que comme prédicteur du recours, jamais comme besoin.

\subsection{Sensibilité : la politique se distingue-t-elle du choix de l'estimateur ?}\label{bes:sensibilite}
Protocole gelé. En 2050, on compare l'effet d'une hausse de 20\,\% des quotas de formation à l'étendue de $C$ entre
estimateurs. Le rapport vaut 1,91\mesure{} (SF), 40\mesure{} (CD) et 1,14\mesure{} (PH) : la politique se
distingue. Il vaut 0,77\mesure{} (MK) et 0,34\mesure{} (IDE) : elle ne se distingue pas. L'étendue vient presque
toute de M0, réfuté par P-B1 ; sans lui, elle reste sous 0,004 en log\mesure{}. M0 est donc retiré de la fourchette
affichée, et le verdict gelé reste publié.

\begin{figure}[h]\centering
\includegraphics[width=0.85\linewidth]{../besoin/figures/c2050_estimateurs.pdf}
\caption{Indice de couverture national en 2050 (2026 = 1), par estimateur et par profession (grille de l'axe~5,
trajectoire centrale du jumeau, Omphale central)\mesure{}. M2a n'est pas défini pour les pharmaciens.}
\label{bes:fig:c2050}
\end{figure}

\paragraph{Le vieillissement change le signe.} Pour les infirmiers, l'indice national 2050 vaut 1,14 si la demande
suit la population (M0), mais 0,83 avec M1\mesure{} (figure~\ref{bes:fig:c2050}). L'offre projetée croît plus vite que
la population, mais moins vite que la demande d'une population qui vieillit. Au département, la médiane de
$C_1/C_0-1$ vaut $-28{,}4$\,\% pour les infirmiers et $-14{,}1$\,\% pour les kinés\mesure{}.

\subsection{Transferts de tâches}\label{bes:transferts}
Un acte passe de la profession A à la profession B à partir de la date d'autorisation $t_0$ :
\[
B'_A=B_A(1-s(t)\,a),\qquad B'_B=B_B+s(t)\,a\,k\,B_A+m(t),\qquad s(t)=s_\infty\bigl(1-e^{-(t-t_0)/\tau}\bigr),
\]
où $a$ est la part du temps de A consacrée à l'acte, $k$ le rapport des durées chez B et chez A, et $m$ l'expansion
du recours. Les effectifs ne changent pas. La médecine générale et les spécialistes sont admis comme
\emph{source} seulement : on rapporte le temps libéré, jamais un indice de couverture pour eux (règle O2).
L'identité de conservation tient à $5{,}7\times10^{-10}$ ETP\mesure{} sur \num{13947} ETP déplacés.

\paragraph{Ce qui est mesurable.} Sur 25 paramètres de transfert et de besoin, 8 sont sourcés, 2 en partie, 1
supposé et 14 non mesurables en données ouvertes\mesure{}. Seule la vaccination antigrippale en officine (L23) a $a$
et $s(t)$ mesurés : $s_\infty=0{,}74\pm0{,}07$\derive{}, $\tau=2{,}4\pm0{,}5$\derive{} ans
(figure~\ref{bes:fig:st}) ; $k$ reste supposé dans $[0{,}25\,;\,1]$\suppose{}.

\begin{figure}[h]\centering
\includegraphics[width=0.55\linewidth]{../besoin/figures/s_t_vaccination.pdf}
\caption{Part des vaccinations antigrippales faites en officine (actes IQVIA rapportés aux doses J07BB
d'Open Medic)\mesure{} et ajustement de $s(t)$ sur 2021--2024\derive{}.}
\label{bes:fig:st}
\end{figure}

\paragraph{Effet servi, après l'erratum P1-4.} La part officinale atteignait déjà 0,687 en 2025--26\mesure{}. Seule
la part au-delà de celle atteinte en 2026 agit : L23 déplace l'indice national 2050 des pharmaciens de
0,03\,\%\mesure{} (1,0126 contre 1,0123) et enlève environ \num{104000} vaccinations aux généralistes en
2050\mesure{}. La première version rejouait la montée en charge depuis 2019 et surestimait l'effet d'un facteur 19
(erratum au registre du 2026-11-15).

\subsection{Ce qui est servi dans l'outil}\label{bes:servi}
\begin{itemize}
  \item \textbf{Indice :} M1 pour les cinq professions, M3 pour les infirmiers ; l'effet du vieillissement est
        montré à part. La fourchette affichée est celle de l'offre ; une fourchette entre estimateurs n'est
        montrée que là où elle existe.
  \item \textbf{Carte :} vue \enquote{Couverture de la demande de ville} au département ; le détail de la méthode est
        replié pour l'administration (chapitre~\ref{sec:produit-ref}).
  \item \textbf{Leviers :} la vaccination en officine est un levier branché de bout en bout, qui change la
        couverture des pharmaciens et non leur nombre. Les autres transferts restent hypothétiques (paramètres
        supposés, signalés comme tels).
  \item \textbf{Assistant :} il appelle l'outil \texttt{couverture} de la passerelle, y compris pour l'effectivité
        d'un transfert (O3).
\end{itemize}

\subsection{Limites}\label{bes:limites}
\begin{itemize}
  \item On mesure la demande de ville privée, pas le besoin de la population ; la part servie par l'hôpital public
        n'est pas mesurée.
  \item L'offre privée du jumeau inclut des salariés (64,5\,\% pour les IDE en 2023\mesure{}) dont les soins en clinique
        ne figurent pas dans le recours de ville : sans effet sur $C$ si cette part est constante, pas sur les
        niveaux.
  \item Le recours n'existe qu'à la région (DAMIR) ou en file active au département : rien ne valide l'indice sous le
        département.
  \item Dentistes retenus de justesse par P-B1 ; sages-femmes perdantes au backtest ; seuils des protocoles non
        étalonnés ; 14 paramètres de transfert sur 25 non mesurables.
\end{itemize}

\paragraph{Reproduire.} \texttt{PYTHONPATH=v2 uv run python -m besoin.tests} (15/15\mesure{}) ; données
\texttt{v2/besoin/donnees/} et leur manifeste ; protocoles gelés et jetons dans \texttt{v2/validation/preregistered/}
et \texttt{horodatage/} ; passages \texttt{v2/besoin/validation/pb*.py}, graine 20261115.

\endgroup
